% !PRINTED 523
\weberpage{523}

Now by § 72, (27), (28) $M$ satisfies a transformation equation which, when
\[
\begin{aligned}
\chi &= M^5 + 5M^4 + 15M^3 + 25M^2 + 25M\\
&= M^3\left[\left(M + \frac{5}{M}\right)^2 + 5\left(M + \frac{5}{M}\right) + 5\right]
\end{aligned}
\]
is set, takes the shape:
\weq{4}{j(\omega) = \gamma_2(\omega)^3 = \frac{(\chi^2 + 10\chi + 5)^3}{\chi}.}
For $m = 19$ one obtains from this the value already found above (§ 125)
\[\gamma_2\left(\frac{-3+\sqrt{-19}}{2}\right) = -96,\]
and for the three remaining values of $m$ one obtains
\weq{5}{\begin{array}{l}
j\left(\dfrac{-7+\sqrt{-51}}{2}\right) = -2^{14}\cdot 27\left(6263 + 1519\sqrt{17}\right),\\[6pt]
\gamma_2\left(\dfrac{-3+\sqrt{-91}}{2}\right) = -48\left(227 + 63\sqrt{13}\right),\\[6pt]
j\left(\dfrac{-1+\sqrt{-99}}{2}\right) = -2^{12}\left(4\,591\,804\,316 + 799\,330\,532\sqrt{33}\right).
\end{array}}
From this one can arrive at cubic equations for
\[f\left(\sqrt{-51}\right)^{24}, \quad f\left(\sqrt{-91}\right)^{8}, \quad f\left(\sqrt{-99}\right)^{24}\]
by making use of the formulae (§ 34, § 54):
\[f_2\left(\frac{-r+\sqrt{-m}}{2}\right) = \frac{\sqrt{2}\, e^{-\frac{r\pi i}{24}}}{f(\sqrt{-m})},\]
\[\gamma_2\left(\frac{-r+\sqrt{-m}}{2}\right) e^{-\frac{r\pi i}{3}} = \frac{f(\sqrt{-m})^{24} - 16^2}{f(\sqrt{-m})^{16}},\]
in which $r = 7, 3, 1$; $m = 51, 91, 99$ are to be set. If, then, one sets for $r = 3$:
\[f(\sqrt{-m}) = x,\]
one obtains:
\weq{6}{x^{24} + \gamma_2(\omega) x^{16} - 16^2 = 0,}
and for $r = 7, 1$:
\[f(\sqrt{-m})^3 = 2x,\]
for these two cases:
\weq{7}{x^{24} - \left[3 - 2^{-8} j(\omega)\right] x^{16} + 3x^8 - 1 = 0,}
%CONT>>
